% Section 4.6, translated from sv/chapters/04/exercises.tex.

\section{Exercises}
\label{c4:sec:uppgifter}

\begin{aside}
\textbf{How to work with the exercises.} Parts~1 and~2 are done during the lesson: first on your
own, a few minutes per exercise, then together. Part~3 is extra exercises for further practice
after the lesson. Most of them can be done in your head or with pencil and paper.

\facitnotis{L04}
\end{aside}

% ------------------------------------------------------------------------------------------
\uppgiftsdel{Part 1}{Review exercises}

\uppgift{1.1}{Ohm's law: find the voltage}
A circuit has $R = 47\,\Omega$ and $I = 200\,\text{mA} = 0.2\,\text{A}$.
\begin{parts}
  \item Calculate the voltage $U$.
  \item The power in the resistor is given by $P = U \cdot I$. Calculate the power.
\end{parts}

\uppgift{1.2}{A linear equation in an RC circuit}
The time constant $\tau$ of an RC circuit is given by $\tau = R \cdot C$. In the circuit,
$\tau = 20\,\text{ms}$ and $C = 10\,\mu\text{F}$.
\begin{parts}
  \item Set up the equation and solve for $R$.
  \item Give $R$ in $\text{k}\Omega$.
\end{parts}

\uppgift{1.3}{Inequality}
An amplifier delivers the output voltage $U_{\text{out}} = 5 \cdot U_{\text{in}}$. The output
voltage must not exceed $15\,\text{V}$. Set up and solve an inequality for $U_{\text{in}}$.

% ------------------------------------------------------------------------------------------
\uppgiftsdel{Part 2}{New material}

\uppgift{2.1}{An algebraic system of equations}
Solve the system of equations
\[
  \begin{cases}
    3x + y = 11 \\
    x - 2y = -2
  \end{cases}
\]

\uppgift{2.2}{Currents from Kirchhoff's current law}
At a node, KCL holds: the sum of the currents in equals the sum of the currents out. The currents
$I_1$ and $I_2$ flow into a node, and the current $I_3 = 5\,\text{A}$ flows out. We also know that
$I_1 = 2 I_2$.
\begin{parts}
  \item Set up a system of equations for $I_1$ and $I_2$.
  \item Solve the system and give the values of the currents.
\end{parts}

\uppgift{2.3}{Voltage and current via KVL}
A circuit consists of a voltage source $E = 12\,\text{V}$ and two resistors $R_1$ and $R_2$ in
series. KVL gives $E = U_1 + U_2$, that is, $U_1 + U_2 = 12$. The voltage drop across $R_2$ is
also three times that across $R_1$: $U_2 = 3 U_1$.
\begin{parts}
  \item Set up the system of equations.
  \item Solve the system and calculate $U_1$ and $U_2$.
  \item The current is $I = 2\,\text{A}$. Calculate $R_1$ and $R_2$ with Ohm's law.
\end{parts}

% ------------------------------------------------------------------------------------------
\uppgiftsdel{Part 3}{Extra exercises}
The exercises below are extra practice, best done on your own after the lesson.

\uppgift{3.1}{The substitution method}
Solve the systems by the substitution method.
\begin{delar}(3)
  \task $\begin{cases} y = 2x + 1 \\ 3x + y = 11 \end{cases}$
  \task $\begin{cases} x = 3y \\ x + 2y = 15 \end{cases}$
  \task $\begin{cases} x + y = 10 \\ x - y = 4 \end{cases}$
\end{delar}

\uppgift{3.2}{The elimination method}
Solve the systems by the elimination method.
\begin{delar}(4)
  \task $\begin{cases} 2x + y = 7 \\ x - y = 2 \end{cases}$
  \task $\begin{cases} 3x + 2y = 16 \\ x - 2y = 0 \end{cases}$
  \task $\begin{cases} 5x + 3y = 21 \\ 2x - y = 4 \end{cases}$
  \task $\begin{cases} 3x + 4y = 10 \\ 2x - 3y = 1 \end{cases}$
\end{delar}

\uppgift{3.3}{Number of solutions}
Decide whether the system has exactly one solution, no solution or infinitely many. Justify your
answer by how the lines lie.
\begin{delar}[column-sep=0.4em](4)
  \task $\begin{cases} x + y = 4 \\ 2x + 2y = 8 \end{cases}$
  \task $\begin{cases} x + y = 4 \\ x + y = 6 \end{cases}$
  \task $\begin{cases} x + y = 4 \\ x - y = 0 \end{cases}$
  \task $\begin{cases} 2x - 4y = 6 \\ -x + 2y = -3 \end{cases}$
\end{delar}

\uppgift{3.4}{Graphical interpretation}
Consider the system
\[
  \begin{cases}
    y = 2x - 1 \\
    y = -x + 5
  \end{cases}
\]
\begin{parts}
  \item Make a table of values for both lines with $x = 0, 1, 2, 3, 4$.
  \item At what $x$ do the two expressions give the same $y$? Read the point of intersection from
    the table.
  \item Solve the system algebraically and check your reading.
  \item What would the table look like if the system had no solution?
\end{parts}

\uppgift{3.5}{Systems with fractions and decimals}
Solve the systems.
\begin{delar}(2)
  \task $\begin{cases} 0.5x + y = 4 \\ x - y = 2 \end{cases}$
  \task $\begin{cases} \dfrac{x}{2} + \dfrac{y}{3} = 4 \\ x - y = 3 \end{cases}$
\end{delar}

\begin{aside}
\note[Tip] Clear the denominators before you choose a method.
\end{aside}

\uppgift{3.6}{KCL at a node}
The currents $I_1$ and $I_2$ flow into a node, while $I_3 = 12\,\text{A}$ flows out. A measurement
also shows that $I_1$ is $4\,\text{A}$ greater than $I_2$.
\begin{parts}
  \item Set up a system of equations for $I_1$ and $I_2$.
  \item Solve the system.
  \item Check the solution in both equations.
\end{parts}

\uppgift{3.7}{Current division in a parallel connection}
Two resistors $R_1$ and $R_2$ are connected in parallel across the voltage $U = 24\,\text{V}$. The
total current is $I = 3\,\text{A}$, and the branch current through $R_1$ is twice that through
$R_2$.
\begin{parts}
  \item Set up a system of equations for the branch currents $I_1$ and $I_2$.
  \item Solve the system.
  \item Calculate $R_1$ and $R_2$ with Ohm's law.
  \item Check that the parallel resistance $R_1 \parallell R_2$ agrees with $\dfrac{U}{I}$.
\end{parts}

\uppgift{3.8}{KVL in a series circuit}
A series circuit with two resistors is supplied with $E = 15\,\text{V}$. The voltage drops satisfy
KVL, and $U_1$ is $3\,\text{V}$ greater than $U_2$.
\begin{parts}
  \item Set up the system of equations.
  \item Solve the system.
  \item The current is $I = 0.3\,\text{A}$. Calculate $R_1$ and $R_2$.
  \item Check that $R_1 + R_2 = \dfrac{E}{I}$.
\end{parts}

\uppgift{3.9}{Operating point: source and load}
A voltage source with internal resistance gives the terminal voltage $U = 12 - 2I$, where $I$ is
the current in amperes. The source is loaded with a $4\,\Omega$ resistor, which by Ohm's law gives
$U = 4I$.
\begin{parts}
  \item Set up the system of equations for $U$ and $I$.
  \item Solve the system.
  \item Check the solution in both equations.
  \item What is the point called where the characteristics of the source and the load intersect?
\end{parts}

\uppgift{3.10}{Designing a voltage divider}
Two resistors are connected in series, and the total resistance is $R_1 + R_2 = 900\,\Omega$. When
the pair is supplied with $U = 9\,\text{V}$, the voltage across $R_1$ is $6\,\text{V}$.
\begin{parts}
  \item Use the voltage division $\dfrac{U_1}{U} = \dfrac{R_1}{R_1 + R_2}$ to get a second
    equation.
  \item Solve the system and give $R_1$ and $R_2$.
  \item Check your answer with the voltage divider formula.
\end{parts}

\uppgift{3.11}{A system with three unknowns}
In a circuit, the following relationships hold between three branch currents (in amperes):
\[
  \begin{cases}
    I_1 + I_2 + I_3 = 9 \\
    I_1 - I_2 = 1 \\
    I_3 = 4
  \end{cases}
\]
\begin{parts}
  \item Solve the system.
  \item Check the solution in all three equations.
  \item Why is the third equation especially convenient to start with?
\end{parts}

\uppgift{3.12}{Setting up the system from a text}
A circuit board is built with two types of resistor, A and B. Five resistors of type A together
with three of type B give $4\,390\,\Omega$ in series. Two of type A together with four of type B
give $3\,660\,\Omega$.
\begin{parts}
  \item Introduce the notation $R_A$ and $R_B$ and set up the system of equations.
  \item Solve the system.
  \item Check your answer in both equations.
\end{parts}

\uppgift{3.13}{Checking a proposed solution}
Decide, \textbf{without} solving the system, whether the given pair of numbers is a solution.
\begin{parts}
  \item $(x, y) = (2, 3)$ for $x + y = 5$ and $2x - y = 1$.
  \item $(x, y) = (1, 4)$ for $3x + y = 7$ and $x - y = -3$.
  \item $(x, y) = (5, 1)$ for $x - 2y = 3$ and $x + y = 7$.
  \item Why is it not enough to check in only one of the equations?
\end{parts}

\uppgift{3.14}{Mixed systems}
Solve by the method you think suits best.
\begin{delar}(3)
  \task $\begin{cases} 4x - y = 5 \\ 2x + 3y = 13 \end{cases}$
  \task $\begin{cases} x + 3y = 1 \\ 2x - y = 9 \end{cases}$
  \task $\begin{cases} 6x + 2y = 10 \\ 3x - y = 7 \end{cases}$
  \task* Briefly justify why you chose each method.
\end{delar}
